\chapter{Complexity — Shannon Probabilistic}
%%% generated by the blueprint pipeline from TCSlib.Complexity.CircuitComplexity.ShannonProbabilistic
\section{Overview}

Declarations of the Lean module \texttt{TCSlib.Complexity.CircuitComplexity.ShannonProbabilistic}.

\section{Declarations}

\begin{theorem}[Counting fan-in-two circuits]
\lean{BoolCircuit.count_le_pow}
\label{BoolCircuit.count_le_pow}
\leanok
\uses{BoolCircuit.card_computable_dag_le, BoolCircuit.computableDAG}
\difficulty{4}
Let $n,S$ be natural numbers with $1 \le n$, $1 \le S$, and $10nS \le 2^n$. The number
of Boolean functions on $n$ inputs computed by well-formed fan-in-two directed acyclic
circuits of size at most $S$ is at most $2^{n\cdot 2(S+1)}$, and $20n(S+1) \le 9\cdot
2^n$.
\end{theorem}

\begin{theorem}[Shannon fraction bound]
\lean{BoolCircuit.computableDAG_shannon_fraction_le}
\label{BoolCircuit.computableDAG_shannon_fraction_le}
\leanok
\uses{BoolCircuit.card_computable_dag_le, BoolCircuit.computableDAG, BoolCircuit.count_le_pow}
\difficulty{5}
For every natural number $n$, the fraction of all Boolean functions on $n$ input bits
that are computed by a fan-in-two circuit of size at most the natural-number quotient
$2^n/(10n)$ is at most $2^{-2^n/10}$.
\end{theorem}

\begin{theorem}[Number of Boolean functions]
\lean{BoolCircuit.card_boolFun}
\label{BoolCircuit.card_boolFun}
\leanok
\difficulty{1}
For every natural number $n$, the set of Boolean functions on assignments to $n$ Boolean
variables has cardinality $2^{2^n}$.
\end{theorem}

\begin{theorem}[Probability of matching one input]
\lean{BoolCircuit.prob_apply_eq}
\label{BoolCircuit.prob_apply_eq}
\leanok
\difficulty{5}
Let $n$ be a natural number, and let $g$ and $x$ be Boolean functions on the $n$
coordinates. If a Boolean function on these coordinates is chosen uniformly at random,
the probability that its value at $x$ equals $g(x)$ is $1/2$.
\end{theorem}

\begin{theorem}[Probability of matching a circuit at one input]
\lean{BoolCircuit.prob_eval_eq_apply}
\label{BoolCircuit.prob_eval_eq_apply}
\leanok
\uses{BoolCircuit.DAGCircuit, BoolCircuit.DAGCircuit.eval, BoolCircuit.prob_apply_eq}
\difficulty{1}
Let $n$ be a natural number, let $C$ be a directed acyclic circuit with $n$ Boolean
inputs, gates ordered so each gate reads only earlier vertices, and a valid output
vertex, and let $x$ be an input. If a Boolean function on these inputs is chosen
uniformly at random, the probability that its value at $x$ equals the value computed by
$C$ at $x$ is $1/2$.
\end{theorem}

\begin{theorem}[Probability of a specified Boolean function]
\lean{BoolCircuit.prob_forall_eq}
\label{BoolCircuit.prob_forall_eq}
\leanok
\uses{BoolCircuit.card_boolFun}
\difficulty{3}
Let $n$ be a nonnegative integer, and choose uniformly from all Boolean functions on $n$
Boolean inputs. For any specified such function $g$, the probability that the chosen
function equals $g$ at every input is $2^{-2^n}$.
\end{theorem}

\begin{theorem}[Probability of a specified circuit function]
\lean{BoolCircuit.prob_computes}
\label{BoolCircuit.prob_computes}
\leanok
\uses{BoolCircuit.DAGCircuit, BoolCircuit.DAGCircuit.eval, BoolCircuit.prob_forall_eq}
\difficulty{1}
Let $n$ be a natural number and let $C$ be a directed acyclic circuit with $n$ Boolean
inputs, gates ordered so each gate reads only earlier vertices, and a valid output
vertex. When a Boolean function on the $n$ inputs is chosen uniformly at random, the
probability that it agrees with the function computed by $C$ on every input is
$2^{-2^n}$.
\end{theorem}

\begin{theorem}[Product formula for circuit agreement]
\lean{BoolCircuit.prob_computes_eq_prod}
\label{BoolCircuit.prob_computes_eq_prod}
\leanok
\uses{BoolCircuit.DAGCircuit, BoolCircuit.DAGCircuit.eval, BoolCircuit.prob_computes, BoolCircuit.prob_eval_eq_apply}
\difficulty{3}
Let $n$ be a natural number and let $C$ be a directed acyclic circuit with $n$ Boolean
inputs, gates ordered so each gate reads only earlier vertices, and a valid output
vertex. Under the uniform distribution on Boolean functions on $n$ inputs, the
probability that a random function agrees with the function computed by $C$ at every
input equals the product, over all inputs $x$, of the probability that it agrees with
the circuit's value at $x$.
\end{theorem}

\begin{theorem}[Probability of matching a finite family]
\lean{BoolCircuit.prob_exists_eq_le}
\label{BoolCircuit.prob_exists_eq_le}
\leanok
\uses{BoolCircuit.prob_forall_eq}
\difficulty{4}
Let $n$ be a nonnegative integer, let $D$ be a finite set of indices, and let $e$ assign
to each index a Boolean function on $
$ Boolean inputs. Choose a Boolean function uniformly from all such functions. The
probability that the chosen function equals $e(d)$ for at least one $d\in D$ is at most
$|D|/2^{2^n}$.
\end{theorem}

\begin{theorem}[Probability bound for small circuits]
\lean{BoolCircuit.prob_computableDAG_le}
\label{BoolCircuit.prob_computableDAG_le}
\leanok
\uses{BoolCircuit.computableDAG, BoolCircuit.prob_exists_eq_le}
\difficulty{4}
For natural numbers $n$ and $S$, consider the set of Boolean functions on $n$ inputs
computable by a well-formed, fan-in-at-most-two acyclic circuit with AND, OR, and NOT
gates, whose gates read only earlier vertices and whose output is a circuit vertex.
Measure circuit size by the number of inputs plus the number of gates, and require this
size to be at most $S$. Under the uniform distribution on all Boolean functions on $n$
inputs, the probability of choosing a function in this set is at most its cardinality
multiplied by $2^{-2^n}$.
\end{theorem}

\begin{theorem}[Counting bound for small circuits]
\lean{BoolCircuit.prob_computableDAG_le_count}
\label{BoolCircuit.prob_computableDAG_le_count}
\leanok
\uses{BoolCircuit.card_computable_dag_le, BoolCircuit.computableDAG, BoolCircuit.prob_computableDAG_le}
\difficulty{2}
For natural numbers $n$ and $S$, consider the set of Boolean functions on $n$ inputs
computable by a well-formed, fan-in-at-most-two acyclic circuit with AND, OR, and NOT
gates, whose gates read only earlier vertices and whose output is a circuit vertex.
Measure circuit size by the number of inputs plus the number of gates, and require this
size to be at most $S$. Under the uniform distribution on all Boolean functions on $n$
inputs, the probability of choosing a function in this set is at most $(S+1)(3(S+1)^2)^S
S\,2^{-2^n}$.
\end{theorem}

\begin{theorem}[Shannon bound for small circuits]
\lean{BoolCircuit.prob_computableDAG_shannon_le}
\label{BoolCircuit.prob_computableDAG_shannon_le}
\leanok
\uses{BoolCircuit.computableDAG, BoolCircuit.computableDAG_shannon_fraction_le, BoolCircuit.prob_computableDAG_le}
\difficulty{4}
For every natural number $n$, consider the set of Boolean functions on $n$ inputs
computable by a well-formed, fan-in-at-most-two acyclic circuit with AND, OR, and NOT
gates, whose gates read only earlier vertices and whose output is a circuit vertex, and
whose size (number of inputs plus gates) is at most $\left\lfloor
2^n/(10n)\right\rfloor$. Under the uniform distribution on all Boolean functions on $n$
inputs, the probability of choosing a function in this set is at most $2^{-2^n/10}$.
\end{theorem}

\begin{theorem}[Real cutoff for computable circuits]
\lean{BoolCircuit.computableDAG_real_cutoff}
\label{BoolCircuit.computableDAG_real_cutoff}
\leanok
\uses{BoolCircuit.DAGCircuit, BoolCircuit.DAGCircuit.IsFaninTwo, BoolCircuit.DAGCircuit.eval, BoolCircuit.DAGCircuit.size, BoolCircuit.computableDAG}
\difficulty{5}
For every natural number $n$, the set of Boolean functions on $n$ input bits computed by
a fan-in-two circuit of size at most $2^n/(10n)$, with the size viewed as a real number,
equals the set of functions computed by such a circuit of size at most the
natural-number quotient $2^n/(10n)$.
\end{theorem}
